\documentclass[../DoAn.tex]{subfiles}
\begin{document}

Đồ án đã xây dựng một hệ thống hỗ trợ phát hiện và phân tích phần mềm
độc hại trên nền tảng Windows, kết hợp phân tích tĩnh bằng học máy với
phân tích động trong môi trường sandbox. Hệ thống vận hành tự động từ
khâu nhận tệp PE đầu vào đến khâu xuất báo cáo, không yêu cầu can thiệp
thủ công, đáp ứng các yêu cầu chức năng và phi chức năng đặt ra tại
Chương~\ref{chapter:Design}. Điểm khác biệt chính so với các công cụ
hiện có là khai thác ETW của .NET Runtime để quan sát hành vi nạp và
sinh mã .NET trong bộ nhớ, đồng thời hợp nhất tín hiệu tĩnh và động để
bù trừ cho nhau.

Giai đoạn tĩnh so sánh ba thuật toán Random Forest, XGBoost và LightGBM
trên cùng tập dữ liệu; LightGBM đạt AUC-ROC cao nhất là 0{,}9993 và
được chọn triển khai. Khảo sát số chiều cho thấy hiệu năng phân loại
bão hòa từ khoảng 64 đặc trưng, tức tín hiệu tập trung ở một số ít đặc
trưng; nhờ đó mô hình rút gọn còn 384 chiều vẫn giữ F1 ở mức 0{,}9615
trên tập EMBER 2024. Ngưỡng phân loại được đặt ở 0{,}33 theo định hướng
ưu tiên giảm bỏ sót mã độc.

Giai đoạn động thu thập sự kiện từ Sysmon và SilkETW, sau đó chuẩn hóa,
dựng cây tiến trình, xây dựng chuỗi hành vi, trích xuất chỉ số xâm phạm
và ánh xạ sang khung MITRE ATT\&CK. Trên 205~mẫu mã độc thuộc bảy nhóm
hành vi (36~họ), 72\% mẫu sinh bằng chứng hành vi cụ thể và hệ thống nhận
diện đúng 204~mẫu. Kết quả của cả hai giai đoạn được tổng hợp thành một
báo cáo JSON thống nhất và hiển thị qua giao diện web.

Quy mô đánh giá động giới hạn ở 305~mẫu (205 mã độc và 100 lành tính)
do chi phí thực thi trong sandbox, nên kết quả mang tính minh chứng hơn
là thống kê trên tập lớn.
Khoảng 17\% mẫu mã độc (34 trong 205) không bộc lộ hành vi trong cửa sổ
quan sát do sử
dụng kỹ thuật né tránh môi trường ảo hóa, vốn nằm ngoài phạm vi xử lý
của đồ án.

Giai đoạn tĩnh vẫn sinh báo động giả trên một số phần mềm hợp pháp được
đóng gói nặng, do cấu trúc tệp gần với mã độc. Cơ chế kết hợp kết quả
hai giai đoạn dựa trên tập quy tắc cố định, chưa học tự động từ dữ liệu
nên còn trường hợp quy tắc tiêm mã kích hoạt sai. Phạm vi đồ án cũng
giới hạn ở tệp PE trên Windows, chưa bao gồm tài liệu chứa macro hay
tệp kịch bản.

Hướng phát triển trước mắt là nâng cao khả năng quan sát đối với mẫu né
tránh sandbox, thông qua giả lập tương tác người dùng, kéo dài thời
gian quan sát và che giấu dấu hiệu môi trường ảo hóa. Cơ chế tổng hợp
kết quả có thể được thay bằng một mô hình học máy kết hợp đặc trưng
tĩnh và động, giảm phụ thuộc vào tập quy tắc thủ công.

Về dữ liệu, quy mô đánh giá động có thể mở rộng bằng cách tự động hóa
việc thực thi song song trên nhiều máy ảo. Về phạm vi, hệ thống có thể
được mở rộng để phân tích thêm tài liệu Office chứa macro và tệp kịch
bản, đồng thời bổ sung phân tích bộ nhớ và lưu lượng mạng nhằm làm giàu
bằng chứng hành vi.

\end{document}
